\lecture{9}{2026-10-08}{Picking the group (and the representation) for Diffie-Hellman}

How to find \(p\) to do Diffie-Hellman in \(\mathbb{Z}_p\)?

The first idea is to pick \(n\) with \(3072\) bits until \(n\) prime.
The problem is that the order of \(\mathbb{Z}_n^*\) is \(n-1\), which is not prime, and there's no guarantee of large prime factors.

A second idea is to pick a \(3072\)-bit prime \(p\) such that \(p = 2q + 1\) for some prime \(q\).
This ensures that the order of \(\mathbb{Z}_p^*\) is \(p-1 = 2q\), which has a large prime factor \(q\).
Then we'd use \(g = x^2 \mod p\) for a random \(x\) to ensure that \(g\) is a generator of the subgroup of order \(q\).
However, the problem is that \(q\) is too big. % What exactly is too big? Cause the computations are still done modulo p.
Here, square-and-multiply has complexity \(O( (\log p)^3 )\).

A third idea is to pick a \(3072\)-bit prime \(p\) of the form \(f q + 1\), for some \(256\)-bit prime \(q\) and some integer \(f\).
This ensures that the order of \(\mathbb{Z}_p^*\) is \(p-1 = f q\), which has a large prime factor \(q\).
Then we'd use \(g = x^f \mod p\) for a random \(x\) to ensure that \(g\) is a generator of the subgroup of order \(q\).
Here, square-and-multiply has complexity \(O( (\log q) (\log p)^2 )\). 

The key here is that all of these are doing some computation in a group isomorphic to a cyclic group of prime order, but since the elements of this group (and therefore the operations of this group) are represented in distinct ways, the complexity of the operations is different.

A fourth idea is to use polynomial rings.
% define irreducible
% 

The final idea is to use elliptic curves, which also give us a group isomorphic to a cyclic group of prime order; that is, we still have the same group, but we're representing the elements of this group in a different way, where the operations have different complexity.

\section*{Elliptic curves}

% \begin{definition}
%     An elliptic curve over a field \(K\) is an equation of the form
%     \[
%         y^2 + a_1 xy + a_3 y = x^3 + a_2 x^2 + a_4 x + a_6
%     \]
%     for some \(a_1, a_2, a_3, a_4, a_6 \in K\).
%     We also identify an elliptic curve with the set of points \((x, y) \in K^2\) satisfying its equation, together with a point \(\infty\) 
% \end{definition}

% If \(\operatorname{char}{K} \neq 2\), we can simplify the equation to the form
% \[
%     y^2 = x^3 + a_2 x^2 + a_4 x + a_6.
% \]
% If moreover \(\operatorname{char}{K} \notin \{2, 3\}\), we can simplify the equation to the form
% \[
%     y^2 = x^3 + ax + b.
% \]

\begin{definition}[Elliptic curve in characteristic not 2 or 3]
    An elliptic curve over a field \(K\) of characteristic not 2 or 3 is an equation of the form
    \[
        y^2 = x^3 + ax + b
    \]
    for some \(a, b \in K\) such that \(4a^3 + 27b^2 \neq 0\).
    We also identify an elliptic curve with the set of points \((x, y) \in K^2\) satisfying its equation, together with a point \(\infty\) at infinity.
\end{definition}

In \(\mathbb{R}\), here are some pictures of elliptic curves.

\begin{figure}[htbp]
    \begin{minipage}{0.5\textwidth}
        \centering
        \begin{tikzpicture}
            \draw[->] (-1.8,0) -- (3.8,0) node[right] {\(x\)};
            \draw[->] (0,-2.8) -- (0,2.8) node[above] {\(y\)};
            \draw[blue, thick, smooth, domain=-1.3247:3, samples=100]
                plot ({1.2*\x}, {0.5*sqrt(\x*\x*\x - \x + 1)});
            \draw[blue, thick, smooth, domain=-1.3247:3, samples=100]
                plot ({1.2*\x}, {-0.5*sqrt(\x*\x*\x - \x + 1)});
        \end{tikzpicture}
    \end{minipage}
    \begin{minipage}{0.5\textwidth}
        \centering
        \begin{tikzpicture}
            \draw[->] (-1.8,0) -- (3.8,0) node[right] {\(x\)};
            \draw[->] (0,-2.8) -- (0,2.8) node[above] {\(y\)};
            \draw[blue, thick, smooth, domain=-1.3247:0.618, samples=100]
                plot ({1.2*\x}, {0.5*sqrt(\x*\x*\x - 2*\x + 1)});
            \draw[blue, thick, smooth, domain=1:3, samples=100]
                plot ({1.2*\x}, {0.5*sqrt(\x*\x*\x - 2*\x + 1)});
            \draw[blue, thick, smooth, domain=-1.3247:0.618, samples=100]
                plot ({1.2*\x}, {-0.5*sqrt(\x*\x*\x - 2*\x + 1)});
            \draw[blue, thick, smooth, domain=1:3, samples=100]
                plot ({1.2*\x}, {-0.5*sqrt(\x*\x*\x - 2*\x + 1)});
        \end{tikzpicture}
    \end{minipage}
\end{figure}


So far, we've defined the set of points on an elliptic curve, but we haven't defined a group operation on this set.

The point at infinity \(\infty\) will be the identity element of the group.
We interpret geometrically that the point at infinity \(\infty\) belongs to every vertical line.

We define the operation (denoted by \(+\)) so that is satisfies the following geometric property:
given any line that intersects with the elliptic curve at three points \(P, Q, R\) (counting multiplicities), we have \(P + Q + R = \infty\).

